\section{正規スライスから回収する素イデアルと葉空間}

\subsection{CGG量子化のイデアルを元の多項式空間で扱う機構}

$N$ が局所冪零、$[N,X]=N$ のCGG積を考える。
$u\in\ker N$、$Xu=au$、$a\ne0$ で、$C[u^{-1}]=K[t]$、
$Nt=1,Xt=t$ がある場合、局所化の内部自己同型は
\[
 u^{-1}\star f\star u=(1+sN)^a f.
\]
両側イデアルに属する $f$ の有限な $N$ 軌道上で、
$((1+sN)^{an}f)_{n=0}^{d}$ から $(N^kf)_{k=0}^{d}$ を補間で回収する。
従って全両側イデアルは $N$ 安定であり、$t$ の最高係数を順に除けば
$I=Q[t]$、$Q\subset K$、$XQ\subset Q$ となる。
Poisson側でも $\{u,f\}=-auNf$ から同じ回収ができる。
さらに
\[
 (Q_1[t])\star(Q_2[t])=(Q_1Q_2)[t]
\]
なので、素性は微分素性へ移り、標数零の通常の素性と一致する。
この補題は、全局所冪零導分が線形であることを必要としない。
固有重量 $a=0$ のchartにはそのまま適用できない。

\subsection{正共鳴の全層と閉包}

$m=2j+r\ge3,r>0$、通常CY値、$s\ne0$ について、
$x$、商での $y$、必要な場合の商での $z$ を順に反転する。
スライスと不変環は
\[
\begin{array}{c|c|c}
 \text{部分}&t&K\\\hline
 x\ne0&y/x&\mathbb C[x^{\pm1},q,J]\\
 x=0,y\ne0&z/(2y)&\mathbb C[y^{\pm1},H_r]\\
 j=0,x=y=0,z\ne0&-w/z^r&\mathbb C[z^{\pm1}]
\end{array}
\]
であり、$H_r=w+(-1)^j y^{2j-1}z^{r+1}/(2r+2)$ である。
残る閉部分では $N=0$ となる。
この順序で、$\operatorname{Spec}S=\operatorname{PSpec}(C,N\wedge X)$ が
同じ多項式の部分空間として成立し、全素イデアルは完全素となる。
集合の対応だけでなく、全包含関係と位相も一致する。

整数重量 $(m+3)X$ と $N$ は二次元可解代数群 $G$ を与える。
階数零の集合は $F=V(x,y)$ ($j\ge1$)、$F=V(x,y,z)$ ($j=0$) であり、
他のシンプレクティック葉は二次元 $G$ 軌道である。
$F$ は $G$ の固定点集合ではない。$F$ の各点は一つの葉なので、葉の飽和は
\[
 \operatorname{Sat}_\pi(T)=G\cdot(T\setminus F)\cup(T\cap F)
\]
としなければならない。
葉閉包のイデアルを取る写像 $\kappa$ は商写像であり、任意の点集合 $T$ に対し
\[
 \overline{\kappa(T)}=\kappa(\overline{\operatorname{Sat}_\pi(T)})
\]
が成立する。原始スペクトルの全域的な位相はこの葉空間で決定した。
中心パラメータだけを使う閉集合の列挙と、部分解消上の代数層は別問題である。

原始スペクトルの各層の内部は、$n=(m+r+3)/\gcd(m+r+3,r)$ として
$A_{n-1}$ 曲面、$\mathbb A^1$、$j=0$ の場合だけ現れる一点、および $F$ である。
非文字表現の原始商は、該当する正規元を反転すると全て $\mathcal D(\mathbb G_m)$ となる。
特に $r\ge2$ で全中心が $\mathbb C$ であっても、原始スペクトルには二次元の開層がある。

\subsection{Dixmier--Moeglin同値に必要な最小の分数体情報}

局所素商の係数体 $K_0$ は整数重量の固有元で生成され、非零重量の可逆元 $u$ を持つ。
$u$ の重量を $a\ne0$、固有元 $f$ の重量を $b$ とすると、
$f^a/u^b$ は不変である。従って
\[
 \operatorname{trdeg}K_0^X=\operatorname{trdeg}K_0-1.
\]
この不変体が分数斜体の中心に入ることだけで、rational性から $\dim K/Q=1$ を得る。
分数斜体の全中心を先に計算する必要はない。
可逆固有元を持つ既約曲線では、全 $\mathbb G_m$ 軌道が開なので一軌道に限られ、
局所量子商は $\mathcal D(\mathbb G_m)$ となる。
これと正規局所化補題を合わせて、量子・Poisson双方のDixmier--Moeglin同値を証明できる。

\subsection{次に使える方向}

全域的な葉空間が得られたため、$r=1$ では $D$ 型中心の部分解消と
原始スペクトルの関係を、集合の対応より強い対象で記述することが次の候補となる。
中心の原点ファイバーの冪零厚みはZariski位相からは失われるので、
代数層を構成するなら冪零指数と両側加群の貼合せを同時に保持する必要がある。
$r=0$ や非CY値へ広げる場合は、ゼロ固有重量のchartを先に特定し、
今回のイデアル補間が成立する範囲を個別に判定する。
